We define recoverable urban resilience as the fraction of cumulative functional loss that can be reduced by optimized intervention relative to a no-additional-intervention baseline. The concept links three layers: underlying functional deficit, locally experienced deficit, and access-weighted resident loss.

The data-mining stage evaluates whether empirical mobility, rainfall, demand, and network data contain the ingredients needed to estimate those layers before the counterfactual optimization model is introduced.
